\documentclass[a4paper]{article}
% Español (México) con XeLaTeX: fontspec para fuentes Unicode, polyglossia
% para la división silábica y los nombres del documento en la variante mexicana.
\usepackage{fontspec}
\usepackage{polyglossia}
\setdefaultlanguage[variant=mexican]{spanish}
% Los nombres chinos van como «pinyin (original)»: xeCJK da a los caracteres
% Han su propia fuente; sin ella XeLaTeX los omite con un aviso.
% FandolSong no cubre algunos caracteres de nombres (U+73FA); WenQuanYi Zen Hei sí.
\usepackage[AutoFallBack=true]{xeCJK}
\setCJKmainfont{FandolSong-Regular.otf}
\setCJKfallbackfamilyfont{\CJKrmdefault}{WenQuanYi Zen Hei}
% polyglossia no traduce los nombres de listings.
\AtBeginDocument{\providecommand{\lstlistingname}{}\renewcommand{\lstlistingname}{Listado}%
  \providecommand{\lstlistlistingname}{}\renewcommand{\lstlistlistingname}{Índice de listados}}
\usepackage{amsmath, amssymb}
\usepackage{graphicx}
\usepackage[margin=2.5cm]{geometry}
\usepackage[most]{tcolorbox}
\usepackage{etoolbox}
\usepackage{listings}
\usepackage{hyperref}
\usepackage{booktabs}
\usepackage{float}
\newtcolorbox{knowledgebox}[1]{enhanced, colback=blue!5!white, colframe=blue!75!black, colbacktitle=blue!75!black, coltitle=white, fonttitle=\bfseries, title={#1}, attach boxed title to top left={yshift=-2mm, xshift=2mm}, boxrule=1pt, sharp corners}
\newtcolorbox{importantbox}[1]{enhanced, colback=yellow!10!white, colframe=yellow!80!black, colbacktitle=yellow!80!black, coltitle=black, fonttitle=\bfseries, title={#1}, sharp corners}
\newtcolorbox{warningbox}[1]{enhanced, colback=red!5!white, colframe=red!75!black, colbacktitle=red!75!black, coltitle=white, fonttitle=\bfseries, title={#1}, sharp corners}
\lstset{language=Python, basicstyle=\ttfamily\small, keywordstyle=\color{blue}, stringstyle=\color{red!60!black}, commentstyle=\color{green!60!black}, breaklines=true, frame=single, numbers=left, numberstyle=\tiny\color{gray}, captionpos=b, extendedchars=true}
\newcommand{\notetitle}{veRL FSDP SFT Trainer explicado en detalle}
\newcommand{\noteauthors}{Elaboradas a partir de materiales públicos del curso}
\newcommand{\notedate}{2025}
\newcommand{\videochannel}{Wudaokou Nashi (五道口纳什)}
\newcommand{\videopublishdate}{2025}
\newcommand{\videoduration}{aproximadamente 25 minutos}
\newcommand{\videourl}{}
\newcommand{\videocoverpath}{cover.jpg}
\newcommand{\repourl}{https://github.com/hqhq1025/ai-course-notes}

% Footer with repo info
\usepackage{fancyhdr}
\pagestyle{fancy}
\fancyhf{}
\fancyfoot[C]{\thepage}
\fancyfoot[R]{\footnotesize\color{gray}\href{\repourl}{AI Course Notes}}
\renewcommand{\headrulewidth}{0pt}
\renewcommand{\footrulewidth}{0pt}

\begin{document}
\begin{titlepage}\centering
{\Large Notas del curso\par}\vspace{1.2cm}{\huge\bfseries \notetitle\par}\vspace{0.8cm}{\large \noteauthors\par}\vspace{0.3cm}{\large \notedate\par}\vspace{1.2cm}
\ifdefempty{\videocoverpath}{}{\includegraphics[width=0.82\textwidth,height=0.45\textheight,keepaspectratio]{\videocoverpath}\par}
\vfill
\begin{tcolorbox}[width=0.9\textwidth, colback=black!2!white, colframe=black!60, sharp corners]
\textbf{Autor o canal del video}: \videochannel\par\textbf{Fecha de publicación}: \videopublishdate\par\textbf{Duración del video}: \videoduration\par
\ifdefempty{\videourl}{}{\textbf{Enlace del video}: \href{\videourl}{\nolinkurl{\videourl}}\par}\par
\textbf{Página del proyecto}: \href{\repourl}{\nolinkurl{\repourl}}\par
\end{tcolorbox}
\end{titlepage}
\tableofcontents\newpage

\section{Introducción}
En esta entrega se presenta el FSDP SFT Trainer de veRL, sencillo y eficiente, y se repasa el proceso completo del entrenamiento de SFT.

\section{Diseño de un buen marco de entrenamiento}
\begin{importantbox}{Dos grandes dimensiones de diseño}
\begin{itemize}
  \item \textbf{Dataset}: abstracción de la carga y el preprocesamiento de datos (SFT Dataset, RLHF Dataset)
  \item \textbf{Trainer}: abstracción del ciclo de entrenamiento (SFT Trainer, PPO Trainer)
\end{itemize}
\end{importantbox}

\section{SFT Dataset}
El núcleo es la función \texttt{get\_item}: convierte los datos prompt-response en la entrada del modelo. Puntos clave:
\begin{itemize}
  \item Solo se calcula la pérdida sobre la parte de response (la parte del prompt se enmascara)
  \item Procesamiento de tokenización: se agregan tokens especiales, padding y attention mask
\end{itemize}

\section{FSDP SFT Trainer}
El código del SFT Trainer de veRL es extremadamente conciso y claro, y es una excelente plantilla para aprender los detalles de implementación del SFT. FSDP (Fully Sharded Data Parallel) implementa el entrenamiento distribuido.

\section{Diferencias de entrenamiento entre SFT y RL}
\begin{itemize}
  \item SFT: aprendizaje supervisado, con un target explícito, no requiere rollout
  \item RL: requiere generar la respuesta mediante rollout en línea, la reward funciona como señal de entrenamiento
\end{itemize}

\subsection{Resumen de la sección}
El SFT Trainer de veRL es una excelente implementación de referencia, con un código conciso y de alta legibilidad.

\newpage
\section{Método de lectura de código: primero ver cómo fluyen los datos hacia el Trainer}
El valor de esta clase está en fijar con claridad las dos abstracciones centrales, Dataset y Trainer. Al leer código de un framework de verdad, lo más útil de seguir primero es: cómo pasa una muestra de \texttt{get\_item} a un batch, y de ahí fluye hacia el Trainer para completar la propagación hacia adelante y hacia atrás.

\begin{knowledgebox}{La ruta más corta para leer un framework de entrenamiento}
\begin{itemize}
  \item Primero ver cómo organiza el Dataset las entradas y salidas
  \item Después ver cómo consume el Trainer el batch
  \item Al final ver cómo envuelve la estrategia distribuida el ciclo de entrenamiento
\end{itemize}
\end{knowledgebox}

\subsection{Resumen de la sección}
Para entender un framework de entrenamiento no es necesario empezar por el código distribuido más complejo; ver primero el flujo de datos y la lógica del Trainer suele facilitar más la construcción de una visión de conjunto.

\newpage
\section{Síntesis y ampliación}
\begin{enumerate}
  \item Dataset + Trainer son las dos abstracciones centrales del marco de entrenamiento
  \item SFT solo calcula la pérdida sobre el response
  \item FSDP implementa el entrenamiento distribuido
\end{enumerate}
\end{document}
